% backmatter/exercise_solutions/chapter2/sec1_solutions.tex

\section*{Chapter 2: Communication and Sensing Fundamentals}
\addcontentsline{toc}{section}{Chapter 2: Communication and Sensing Fundamentals}

\subsection*{Section 2.1 Exercises}

\begin{enumerate}[label=\textbf{Exercise 2.\arabic*.}, leftmargin=*]
    \item The wavelength is
    \[
        \lambda=\frac{c}{f_c}=\frac{3\times10^8}{3\times10^9}=0.1\,\mathrm{m}.
    \]
    The free-space path loss is
    \[
        \mathrm{FSPL}_{\mathrm{dB}}
        =
        20\log_{10}\left(\frac{4\pi d}{\lambda}\right)
        =
        20\log_{10}\left(\frac{4\pi\times1000}{0.1}\right)
        \approx 102\,\mathrm{dB}.
    \]
    With unit antenna gains,
    \[
        P_r \approx 0\,\mathrm{dBm}-102\,\mathrm{dB}
        =
        -102\,\mathrm{dBm}.
    \]
    The noise power is
    \[
        P_n=-174+10\log_{10}(10\times10^6)
        =
        -104\,\mathrm{dBm}.
    \]
    Hence
    \[
        \mathrm{SNR}_{\mathrm{dB}}
        =
        P_r-P_n
        \approx 2\,\mathrm{dB}.
    \]
    In linear scale, \(\mathrm{SNR}\approx 10^{2/10}\approx1.59\). The
    capacity is
    \[
        C=B\log_2(1+\mathrm{SNR})
        \approx
        10^7\log_2(1+1.59)
        \approx 13.7\,\mathrm{Mbit/s}.
    \]

    \item Taking \(P_s\) as the reference, the interference power is \(6\) dB
    below \(P_s\), and the noise power is \(10\) dB below \(P_s\). Therefore,
    \[
        \frac{P_i}{P_s}=10^{-6/10}\approx0.251,
        \qquad
        \frac{P_n}{P_s}=10^{-10/10}=0.1.
    \]
    The SINR is
    \[
        \mathrm{SINR}
        =
        \frac{P_s}{P_i+P_n}
        =
        \frac{1}{0.251+0.1}
        \approx2.85,
    \]
    or
    \[
        \mathrm{SINR}_{\mathrm{dB}}\approx10\log_{10}(2.85)\approx4.5\,\mathrm{dB}.
    \]
    The approximate throughput is
    \[
        R_{\mathrm{throughput}}
        \approx R(1-\mathrm{BLER})
        =
        20(1-0.1)
        =
        18\,\mathrm{Mbit/s}.
    \]

    \item The maximum Doppler shift is
    \[
        f_D=\frac{v f_c}{c}
        =
        \frac{30\times28\times10^9}{3\times10^8}
        =
        2800\,\mathrm{Hz}.
    \]
    Thus,
    \[
        T_c\approx\frac{1}{f_D}\approx3.57\times10^{-4}\,\mathrm{s}
        =
        357\,\mu\mathrm{s}.
    \]
    Since \(T_s=20\,\mu\mathrm{s}\ll T_c\), the channel is closer to slow
    fading. The coherence bandwidth is
    \[
        B_c\approx\frac{1}{\tau_{\mathrm{rms}}}
        =
        \frac{1}{0.5\times10^{-6}}
        =
        2\,\mathrm{MHz}.
    \]
    Since \(B=100\,\mathrm{MHz}>B_c\), the channel is closer to
    frequency-selective fading.

    \item The MIMO channel matrix has dimension
    \[
        \mathbf{H}[k]\in\mathbb{C}^{N_r\times N_t}
        =
        \mathbb{C}^{2\times4}.
    \]
    Since \(N_s=2\),
    \[
        \mathbf{s}[k]\in\mathbb{C}^{2\times1},
        \quad
        \mathbf{W}[k]\in\mathbb{C}^{4\times2},
        \quad
        \mathbf{x}[k]\in\mathbb{C}^{4\times1},
        \quad
        \mathbf{y}[k]\in\mathbb{C}^{2\times1}.
    \]
    The maximum number of supportable spatial streams is limited by
    \[
        N_s\leq\operatorname{rank}(\mathbf{H}[k])=2.
    \]
    Therefore, at most two spatial streams can be supported on this subcarrier.

    \item The range is obtained from the round-trip delay:
    \[
        R=\frac{c\tau}{2}
        =
        \frac{3\times10^8\times0.8\times10^{-6}}{2}
        =
        120\,\mathrm{m}.
    \]
    The wavelength is
    \[
        \lambda=\frac{c}{f_c}
        =
        \frac{3\times10^8}{77\times10^9}
        \approx3.90\times10^{-3}\,\mathrm{m}.
    \]
    The radial velocity is
    \[
        v=\frac{\lambda f_D}{2}
        =
        \frac{3.90\times10^{-3}\times3000}{2}
        \approx5.84\,\mathrm{m/s}.
    \]

    \item The range resolution is
    \[
        \Delta R=\frac{c}{2B}
        =
        \frac{3\times10^8}{2\times500\times10^6}
        =
        0.3\,\mathrm{m}.
    \]
    With \(f_c=77\,\mathrm{GHz}\),
    \[
        \lambda=\frac{c}{f_c}
        =
        \frac{3\times10^8}{77\times10^9}
        \approx3.90\,\mathrm{mm}.
    \]
    The velocity resolution is
    \[
        \Delta v\approx\frac{\lambda}{2T_{\mathrm{obs}}}
        =
        \frac{3.90\times10^{-3}}{2\times20\times10^{-3}}
        \approx0.097\,\mathrm{m/s}.
    \]

    \item The phase difference of adjacent array elements is
    \[
        \Delta\phi=\frac{2\pi d\sin\theta}{\lambda}.
    \]
    Since \(d=\lambda/2\),
    \[
        \Delta\phi=\pi\sin\theta.
    \]
    With \(\Delta\phi=\pi/3\),
    \[
        \sin\theta=\frac{1}{3},
        \qquad
        \theta=\sin^{-1}\left(\frac{1}{3}\right)\approx19.5^\circ.
    \]

    \item The Neyman--Pearson threshold satisfies
    \[
        P(T>\gamma|H_0)=0.05.
    \]
    Since \(T|H_0\sim\mathcal{N}(0,1)\),
    \[
        \gamma=\Phi^{-1}(0.95)\approx1.645.
    \]
    Under \(H_1\), \(T\sim\mathcal{N}(3,1)\). Therefore,
    \[
        P_D=P(T>\gamma|H_1)
        =
        1-\Phi(\gamma-3)
        =
        1-\Phi(-1.355)
        =
        \Phi(1.355)
        \approx0.91.
    \]
\end{enumerate}
